% Fichier complété par tools/complete_missing_corrections.py.
% Source : GEO/GEO-01/09_RT_RSG/09_RT_RSG.tex
% Statut : rédigé (3 question(s)).
\begin{corrigebox}[Corrigé rédigé]
\CorrigeQuestion{1}
Le graphe comporte 0 (sol), 1 (roue) relié à 0 par
                    un contact de roulement sans glissement en I, puis 2 relié à 1 par un
                    pivot en A. Pour 09, la glissière 2/1 est ajoutée suivant $\vec i_1$;
                    pour 46, 2/1 est un pivot.
\CorrigeQuestion{2}
On reconstruit la configuration en appliquant les
                rotations dans l'ordre de la chaîne cinématique puis les translations le
                long de leurs axes. Les longueurs des barres restent constantes. Les
                coordonnées finales se calculent avec
                $\vec i'=\cos q\,\vec i+\sin q\,\vec j$; le tracé doit conserver les
                liaisons et le contact sans pénétration.
\CorrigeQuestion{3}
On reconstruit la configuration en appliquant les
                rotations dans l'ordre de la chaîne cinématique puis les translations le
                long de leurs axes. Les longueurs des barres restent constantes. Les
                coordonnées finales se calculent avec
                $\vec i'=\cos q\,\vec i+\sin q\,\vec j$; le tracé doit conserver les
                liaisons et le contact sans pénétration.
\end{corrigebox}
